\documentclass{article}
\usepackage[margin=1in]{geometry}
\usepackage{hyperref}
\usepackage{fancyhdr}
\usepackage{float}
\usepackage{placeins}

\pagestyle{fancy}
\fancyhead{}
\fancyfoot[L]{EagleNav}
\fancyfoot[C]{Software Design Document}
\fancyfoot[R]{\thepage}

\title{Software Design Document (SDD) - EagleNav}
\author{Group 5}
\date{May 2026}


\begin{document}
\maketitle
\clearpage

\tableofcontents
\newpage

\section{Version History}
\begin{table}[h!]
	\centering
	\begin{tabular}{|c|c|}
		\hline
		\textbf{Snapshot} & \textbf{Date} \\
		\hline
		Snapshot 1: Initial version & 5/7/2026 \\
		Snapshot 2: Added features & 5/8/2026 \\
		Snapshot 3: Updated system design & 5/12/2026 \\
		Snapshot 4: Final version & 5/14/2026 \\
		\hline
	\end{tabular}
\end{table}


\section{Introduction}
\subsection{Purpose}
This document describes the software design of the EagleNav system.
EagleNav is an augmented reality navigation platform. EagleNav aims to remove navigation barriers and create a more inclusive campus experience for all students.


\subsection{Intended Audience & Reading Suggestions}
This document is intended for students, developers, and instructors.
\begin{itemize}
	\item Students: sections 2,4 
	\item Developers: section 3,4 
	\item Instructors: sections 2,3,4
\end{itemize}

\subsection{Overview}
EagleNav is a campus navigation system design to provide routing with a focus on accesibility, 
safety, and real-time awareness. The platform integrates multiple data sources including campus map data, routing engines, sensor inputs, and external services. 
The components work together to generate navigation instructs that adapt to user movement.


\section{System Architecture}
The system includes a mobile frontend, backend server, and database. It is layered into cloud infrstructure, on-device controllers, and on-device services.

\subsection{Architecture Breakdown}
The system is divided into the following major layers: 

\subsection{Cloud infrastructure Layer}
The cloud layer is reponsible for storing map data and generating routes.

\begin{itemize}
	\item Campus map data is stored in Cloudflare R2
	\item The Valhalla routing engine computes possible routes.
	\item Updating map data requires uploading a new file and restarting routing service.
\end{itemize}

\subsubsection{Client Controller Layer}
The client layer runs on the user's device and manages navigation experience.
\begin{itemize}
	\item \textbf{Location Controller}: Handles GPS input
	\item \textbf{Routing Controller}: Requests routes ffrom cloud service
	\item \textbf{Guidance Controller}: Manages navigation
	\item \textbf{Navigation Voice Controller}: Handles spoken instructions
	\item \textbf{UI Navigation Controller}: Controls navigation states: idle, navigating, and arrived
\end{itemize}

\subsection{Service Layer}
The service layer provides reusable system functions.

\begin{itemize}
	\item \textbf{Location Service}: Provides GPS updates
	\item \textbf{Valhalla Service}: Sends routing requests and processes responses
	\item \textbf{Compass Service}: Suplies heading informaiton
	\item \textbf{Text-To-Speech Service}: Generates spoken instructions
	\item \textbf{Haptic Feedback Service}: Provides vibration cues for navigation
\end{itemize}

\subsection{Workflow}
\begin{enumerate}
	\item The user selects a destination in the application
	\item The Routing Controller sends a request to Valhalla
	\item The cloud service returns a route
	\item The Guidance Controller begins navigation
	\item The Location Controller updates the user position
	\item If user deviates from route, The Routing Controller triggers a reroute
	\item Updated routes are recieved
	\item The system ends navigation when user arrives/cancels
\end{enumerate}

\subsection{Workflow Diagram}

\begin{figure}[H]
\centering
\resizebox{\textwidth}{!}{
\begin{tabular}{c}
    
    \fbox{User Selects Destination} \\
    $\downarrow$ \\
    \fbox{Navigation Controller Validates Input} \\
    $\downarrow$ \\
    \fbox{Routing Controller \newline Requests Route (Valhalla)} \\
    $\downarrow$ \\
    \fbox{Route Generated} \\
    $\downarrow$ \\
    \fbox{Guidance Controller Activated} \\
    $\downarrow$ \\
    \fbox{Sensor Fusion + Map Matching Updates} \\
    $\downarrow$ \\
    
    \fbox{
    \begin{tabular}{c}
    On Route $\rightarrow$ Continue Guidance \\
    Off Route $\rightarrow$ Trigger Reroute
    \end{tabular}
    } \\
    
    $\downarrow$ \\
    \fbox{UI Updates (On Track / Rerouting / Arrived)} \\
    $\downarrow$ \\
    \fbox{Navigation Ends (Cancel / Arrival)} \\

    \end{tabular}
    }
    \caption{Compact Navigation Workflow}
	\end{figure}
	\floatbarrier

\section{User Interface}
Users interact through a mobile app with AR directions and voice guidance. The following list represents the primary pages that users will interact with:

\begin{itemize}
\item \textbf{Home Screen:} Users will be able to access the navigation screen, settings, and get basic information at a glance.
\item \textbf{Navigation Screen:} Users will be able to navigate the campus using AR + Map Navigation features.
\item \textbf{Events Screen:} Allows users to view and save campus events.
\item \textbf{Settings:} Users will be able to locally change the app's theme, navigation volume, AR settings, etc.
\end{itemize}

\subsection{Screen Components}
\begin{itemize}
	\item \textbf{HomeMap Screen:} Provides a search bar, map view, and menu.
	\item \textbf{Guidance Screen}: Includes a turnbanner, progess bar of navigation, and an end/cancel option.
\end{itemize}

\subsection{User Interface Flow Diagram}

\begin{figure}[H]
\centering
\resizebox{\textwidth}{!}{
\begin{tabular}{c}

\fbox{HomeMap Screen} \\
$\downarrow$ \\
\fbox{Search / Select Destination} \\
$\downarrow$ \\
\fbox{Route Preview Screen} \\
$\downarrow$ \\
\fbox{Start Navigation} \\
$\downarrow$ \\
\fbox{Guidance Screen (Turn-by-Turn Navigation)} \\
$\downarrow$ \\
\fbox{
\begin{tabular}{c}
On Route $\rightarrow$ Continue Navigation \\
Off Route $\rightarrow$ Recalculate Route
\end{tabular}
} \\
$\downarrow$ \\
\fbox{Arrival Screen (Destination Reached)} \\
$\downarrow$ \\
\fbox{End Navigation}

\end{tabular}
}
\caption{User Interface Flow of EagleNav}
\end{figure}
\floatbarrier


\section{Glossary}
\begin{table}[h!]
	\begin{tabular}{|c|c|}
		\hline
		\textbf{Term} &\textbf{Definition}\\
		\hline
		AR & Augmented Reality \\
		UI & User Interface \\
		\hline
	\end{tabular}
\end{table}



\section{References}
\begin{itemize}
\item{\href{https://ascent.cysun.org/project/project/view/248}{Ascent - Project (https://ascent.cysun.org/project/project/view/248)}}
\end{itemize}

\end{document}
